\section{圆角、目标线与条尾标签}
\begin{example}{示例 37：圆角与真实长度}
\gradbar[rounded=2pt]{72}

\gradbar[rounded=10pt]{1}

\gradbar[rounded=10pt]{0}
\end{example}

\key{rounded} 为圆角半径，默认 \key{0pt}。实际半径会被限制在轨道高度和
彩色条宽度允许的范围内。小数值不被放大，零值仍只显示空轨道。

\begin{example}{示例 38：目标线}
\gradbar[target=80]{72}

\gradbar[target=80,
  target color=red!70!black]{92}
\end{example}

\key{target} 使用与输入值相同的单位，不是相对位置。目标线必须位于
\key{min} 到 \key{max} 之间，越界会报错。
默认不绘制目标线；\key{target={}} 可清除局部或继承的目标设置。

\begin{example}{示例 39：标签跟随彩色条末端}
\gradbar[label=end]{15}

\gradbar[label=end]{72}

\gradbar[label=end]{100}
\end{example}

\key{label=end} 把标签放在彩色条末端上方。靠近左右边界时，
会在轨道范围内向内调整标签位置，避免常规数值标签越界。
因此它会增加行高，但不为外部标签槽预留水平空间。

\note{很长的自定义标签可能比轨道本身还宽；此时应增大 \key{width}，
或改用 \key{label=outside} 并设置足够的 \key{label width}。
上方标签不改变正文中的轨道基线。}
